\chapter{Prompt Library\\คลัง Prompt สำหรับ Mentor}

Prompt Library ไม่ใช่ชุดคำสั่งที่ให้คำตอบสำเร็จรูป แต่เป็นจุดเริ่มต้นสำหรับงานที่เกิดซ้ำในการ Mentoring ทุก Prompt ควรถูกปรับตามเป้าหมายของ Mentee ตรวจสอบข้อมูลก่อนใช้ และนำผลลัพธ์กลับเข้าสู่บทสนทนาระหว่างมนุษย์

\learningobjectives{
  \item เลือก Prompt ให้เหมาะกับช่วงของ Mentoring
  \item ปรับ Context และ Format เพื่อให้ผลลัพธ์นำไปใช้ได้
  \item ตรวจและพัฒนาผลจาก AI อย่างเป็นระบบ
}

\section{Portfolio Builder}
\begin{promptbox}
\textbf{Role:} ท่านเป็นบรรณาธิการ Tech Portfolio สำหรับนักวิจัยและนักนวัตกรรม\\
\textbf{Task:} วิเคราะห์ข้อมูลที่ให้และสร้างร่าง Tech Portfolio ตาม 7 องค์ประกอบ\\
\textbf{Context:} [วาง CV, รายการผลงาน และคำอธิบายโครงการ] ใช้เฉพาะข้อมูลนี้ ห้ามเติมข้อเท็จจริง หากข้อมูลไม่พอให้ระบุคำถาม\\
\textbf{Format:} Profile, Skills, Projects, Process, Evidence, Impact, Reflection แต่ละหัวข้อประกอบด้วยร่างสั้น หลักฐานที่รองรับ และข้อมูลที่ยังขาด
\end{promptbox}
\textbf{วิธีใช้:} ให้ Mentee ตรวจทุกข้อกล่าวอ้าง เลือกสิ่งที่สะท้อนเป้าหมายจริง และใช้ช่องว่างเป็นวาระสนทนา ไม่ควรนำร่างไปเผยแพร่ทันที

\section{Gap Analysis}
\begin{promptbox}
\textbf{Role:} ท่านเป็นผู้ช่วย Mentor ที่วิเคราะห์อย่างระมัดระวังและไม่ตัดสินบุคคล\\
\textbf{Task:} ระบุช่องว่าง 3 ข้อที่มีผลต่อเป้าหมาย [ระบุเป้าหมาย]\\
\textbf{Context:} [วาง Portfolio ที่ลดทอนข้อมูลส่วนบุคคล] แยกสิ่งที่มีหลักฐาน สิ่งที่เป็นข้อสันนิษฐาน และสิ่งที่ไม่ทราบ\\
\textbf{Format:} สำหรับแต่ละ Gap ให้ระบุความสำคัญ สัญญาณจากข้อมูล คำถามเปิด 2 ข้อ และการทดลองขนาดเล็ก 1 อย่าง
\end{promptbox}
\textbf{วิธีใช้:} ตรวจว่า Gap สอดคล้องกับนิยามความสำเร็จของ Mentee หรือไม่ แล้วเลือกเพียงหนึ่งเรื่องเพื่อเจาะลึก

\section{Rewrite and Polish}
\begin{promptbox}
\textbf{Role:} ท่านเป็นบรรณาธิการข้อเสนอด้านนวัตกรรมที่ให้ความสำคัญกับความถูกต้อง\\
\textbf{Task:} ปรับข้อความให้ชัด กระชับ และอ่านออกสำหรับ [กลุ่มผู้อ่าน]\\
\textbf{Context:} [วางข้อความ] รักษาความหมายและตัวเลขเดิม ห้ามเพิ่มคำกล่าวอ้าง แยกผลที่เกิดแล้วออกจากผลที่คาดหวัง\\
\textbf{Format:} แสดงฉบับปรับแล้ว ตามด้วยรายการการเปลี่ยนแปลงสำคัญและคำถามที่ต้องยืนยัน
\end{promptbox}
\textbf{วิธีใช้:} เปรียบเทียบฉบับเดิมกับฉบับใหม่ ให้ Mentee อธิบายว่าการปรับใดช่วยหรือบิดเบือนความหมาย

\section{Infographic Brief}
\begin{promptbox}
\textbf{Role:} ท่านเป็น visual storyteller สำหรับผู้อ่านที่ไม่ใช่ผู้เชี่ยวชาญ\\
\textbf{Task:} เปลี่ยน Portfolio เป็น brief สำหรับ infographic หนึ่งหน้า\\
\textbf{Context:} [วาง Portfolio] กลุ่มผู้อ่านคือ [ระบุ] เป้าหมายการสื่อสารคือ [ระบุ]\\
\textbf{Format:} 1) headline ไม่เกิน 12 คำ 2) Profile 3) ปัญหาและวิธีทำงาน 4) Evidence และ Impact ที่ตรวจสอบได้ 5) next step 6) แนวคิดภาพ โดยทำเครื่องหมายข้อมูลที่ยังขาด
\end{promptbox}
\textbf{วิธีใช้:} ตรวจว่าภาพและตัวเลขไม่ทำให้ผลดูใหญ่เกินจริง และให้เจ้าของข้อมูลอนุมัติก่อนเผยแพร่

\section{วงจรปรับ Prompt}
เมื่อผลลัพธ์ยังไม่ดี อย่าเพียงสั่งว่า \enquote{เขียนใหม่ให้ดีขึ้น} ให้ระบุว่าอะไรไม่ตรง เช่น ผู้อ่านไม่ชัด หลักฐานไม่พอ ภาษากล่าวอ้างเกินจริง หรือรูปแบบใช้ต่อไม่ได้ การปรับ Prompt เป็นการทำให้เกณฑ์การคิดของ Mentor ชัดขึ้นด้วย

\begin{mentornote}
เก็บ Prompt ที่ใช้ได้พร้อมตัวอย่าง input ที่ลดทอนข้อมูลและบันทึกว่าเหมาะกับสถานการณ์ใด Prompt Library ที่ดีควรมีประวัติการเรียนรู้ ไม่ใช่เพียงข้อความคำสั่ง
\end{mentornote}

\begin{keytakeaway}
คุณภาพของ Prompt Library อยู่ที่การนำไปใช้ใน workflow ที่มีเป้าหมาย ข้อมูลเหมาะสม การตรวจสอบ และบทสนทนาต่อเนื่อง มิใช่จำนวน Prompt
\end{keytakeaway}

\begin{reflection}
เลือก Prompt หนึ่งชุด ทดลองกับข้อมูลตัวอย่าง แล้วประเมินผลด้วยเกณฑ์: ถูกต้อง มีประโยชน์ ซื่อตรง ใช้งานได้ และเคารพความลับ ปรับ Prompt หนึ่งรอบตามสิ่งที่พบ
\blankline\blankline
\end{reflection}

\chaptersummary{Prompt ทั้งสี่รองรับงานหลักตั้งแต่สร้าง Portfolio วิเคราะห์ Gap ปรับภาษา ไปจนถึงออกแบบการสื่อสาร ทุกชุดต้องเริ่มจากเป้าหมายและจบด้วยการตรวจสอบของ Mentor และ Mentee}
